\section{Mass-depth rigidity audits and matching-layer details (supplement)}
\label{app:mass_rigidity_details}

This appendix collects audit tables and matching-layer bookkeeping details that support the mass-spectrum closure of Section~\ref{sec:mass_spectrum_closure}.
The main text records the closed template and the spectrum table; the material below provides bounded-complexity evidence and quantized matching summaries.

\subsection{Rigidity of the depth formula at bounded coefficient complexity}
\label{subsec:mass_depth_rigidity}

To make the integer depth assignment auditable as a low-complexity closure, we test a bounded-coefficient ansatz for the depth map.
Write the stable-type differences relative to the electron reference $w_e$ as
\[
\Delta V:=V(w)-V(w_e),\qquad \Delta g:=g(w)-g(w_e),\qquad \Delta|w|_1:=|w|_1-|w_e|_1.
\]
Consider the integer ansatz
\begin{equation}\label{eq:rhat_abc_ansatz}
\widehat r(w)=a\,\Delta V+b\,\Delta g+c\,\Delta|w|_1,
\qquad a,b,c\in\ZZ,
\end{equation}
and measure its mismatch on the scheme-stable charged-lepton anchor set $\{\mu,\tau\}$ by the depth deviation in the resolution coordinate $r(\mu)$.
As a diagnostic, we also record the same metrics on an extended fermion set including quark reference masses (scheme-dependent by convention).
Table~\ref{tab:mass_depth_rigidity} records the selected minimizers in the coefficient box $|a|,|b|,|c|\le B$ for $B=1,\dots,20$ under a lexicographic minimization rule: first minimize the leptonic minimax deviation, then the leptonic sum deviation, then the extended-set minimax and sum deviations, and finally coefficient complexity.

\begin{proposition}[Rigidity of the $(2,5,1)$ coefficients at bounded complexity]\label{prop:rhat_rigidity}
In the bounded search domain $|a|,|b|,|c|\le 20$, the unique minimizer of the leptonic objective (as defined above) satisfies
\[
(a,b,c)=(2,5,1).
\]
Moreover, within the tested range $B=1,\dots,20$, the minimizer stabilizes at $B=5$ and remains constant for all $B$ with $5\le B\le 20$.
\end{proposition}
\begin{proof}
This is a finite exhaustive enumeration over integer triples in the stated box with deterministic tie-break rules; the resulting minimizers are listed in Table~\ref{tab:mass_depth_rigidity} and reproduced by the accompanying script.
\end{proof}

\begin{table}[t]
\centering
\small
\setlength{\tabcolsep}{6pt}
\renewcommand{\arraystretch}{1.15}
\begin{tabular}{rrrrrr}
\toprule
$B$ & $(a,b,c)$ & max$_{\mu,\tau}$ & sum$_{\mu,\tau}$ & max$_{\mathrm{ext}}$ & sum$_{\mathrm{ext}}$ \\
\midrule
1 & $(1,-1,1)$ & 5.079526 & 6.024227 & 10.246816 & 24.792808 \\
2 & $(2,2,-2)$ & 3.079526 & 6.134824 & 7.277571 & 31.990963 \\
3 & $(2,3,-1)$ & 2.079526 & 4.134824 & 6.277571 & 25.990963 \\
4 & $(2,4,0)$ & 1.079526 & 2.134824 & 5.277571 & 19.990963 \\
5 & $(2,5,1)$ & 0.079526 & 0.134824 & 4.277571 & 14.795234 \\
6 & $(2,5,1)$ & 0.079526 & 0.134824 & 4.277571 & 14.795234 \\
7 & $(2,5,1)$ & 0.079526 & 0.134824 & 4.277571 & 14.795234 \\
8 & $(2,5,1)$ & 0.079526 & 0.134824 & 4.277571 & 14.795234 \\
9 & $(2,5,1)$ & 0.079526 & 0.134824 & 4.277571 & 14.795234 \\
10 & $(2,5,1)$ & 0.079526 & 0.134824 & 4.277571 & 14.795234 \\
11 & $(2,5,1)$ & 0.079526 & 0.134824 & 4.277571 & 14.795234 \\
12 & $(2,5,1)$ & 0.079526 & 0.134824 & 4.277571 & 14.795234 \\
13 & $(2,5,1)$ & 0.079526 & 0.134824 & 4.277571 & 14.795234 \\
14 & $(2,5,1)$ & 0.079526 & 0.134824 & 4.277571 & 14.795234 \\
15 & $(2,5,1)$ & 0.079526 & 0.134824 & 4.277571 & 14.795234 \\
16 & $(2,5,1)$ & 0.079526 & 0.134824 & 4.277571 & 14.795234 \\
17 & $(2,5,1)$ & 0.079526 & 0.134824 & 4.277571 & 14.795234 \\
18 & $(2,5,1)$ & 0.079526 & 0.134824 & 4.277571 & 14.795234 \\
19 & $(2,5,1)$ & 0.079526 & 0.134824 & 4.277571 & 14.795234 \\
20 & $(2,5,1)$ & 0.079526 & 0.134824 & 4.277571 & 14.795234 \\
\bottomrule%
\end{tabular}
\caption{Bounded-coefficient rigidity search for the depth ansatz \eqref{eq:rhat_abc_ansatz} over $B=1,\dots,20$. The primary anchor set is $\{\mu,\tau\}$; the extended set includes quark reference masses as a diagnostic. Rows are generated by \texttt{scripts/exp\_mass\_depth\_rigidity.py}.}
\label{tab:mass_depth_rigidity}
\end{table}

\paragraph{Leave-one-out stability of the coefficient choice.}
As an additional audit, we repeat the $B=20$ coefficient search while leaving out one anchor at a time from the finite anchor set $\{u,d,s,c,b,t,\mu,\tau\}$ used by the rigidity script, and record whether the selected minimizer changes under the same lexicographic objective.
Table~\ref{tab:mass_depth_leave_one_out} reports the resulting leave-one-out sweep.

\begin{table}[t]
\centering
\small
\setlength{\tabcolsep}{8pt}
\renewcommand{\arraystretch}{1.15}
\begin{tabular}{llrrrrl}
\toprule
leave-out & $(a,b,c)$ & max$_{\mu,\tau}$ & sum$_{\mu,\tau}$ & max$_{\mathrm{ext}}$ & sum$_{\mathrm{ext}}$ & status \\
\midrule
\texttt{none} & $(2,5,1)$ & 0.079526 & 0.134824 & 4.277571 & 14.795234 & SAME \\
\texttt{-u} & $(2,5,1)$ & 0.079526 & 0.134824 & 4.277571 & 11.790788 & SAME \\
\texttt{-d} & $(2,5,1)$ & 0.079526 & 0.134824 & 4.277571 & 14.393098 & SAME \\
\texttt{-s} & $(2,5,1)$ & 0.079526 & 0.134824 & 4.277571 & 13.609571 & SAME \\
\texttt{-c} & $(1,1,6)$ & 0.079526 & 0.134824 & 3.402135 & 8.177093 & DIFF \\
\texttt{-b} & $(2,5,1)$ & 0.079526 & 0.134824 & 4.246816 & 10.517663 & SAME \\
\texttt{-t} & $(2,5,1)$ & 0.079526 & 0.134824 & 4.277571 & 13.251455 & SAME \\
\texttt{-mu} & $(2,5,1)$ & 0.055298 & 0.055298 & 4.277571 & 14.715708 & SAME \\
\texttt{-tau} & $(2,6,1)$ & 0.079526 & 0.079526 & 4.246816 & 11.281053 & DIFF \\
\bottomrule%
\end{tabular}
\caption{Leave-one-out robustness diagnostic for the depth-coefficient search at $B=20$. Each row removes one anchor from the finite anchor set and recomputes the selected minimizer under the same lexicographic objective as in Table~\ref{tab:mass_depth_rigidity}. Rows are generated by \texttt{scripts/exp\_mass\_depth\_leave\_one\_out.py}.}
\label{tab:mass_depth_leave_one_out}
\end{table}

\begin{table}[t]
\centering
\small
\setlength{\tabcolsep}{10pt}
\renewcommand{\arraystretch}{1.15}
\begin{tabular}{rrrrl}
\toprule
tested leave-outs & SAME & DIFF & fraction SAME & DIFF cases \\
\midrule
8 & 6 & 2 & \texttt{-c}, \texttt{-tau} \\
\bottomrule%
\end{tabular}
\caption{Compact summary of the leave-one-out robustness sweep in Table~\ref{tab:mass_depth_leave_one_out}. The baseline is $(a,b,c)=(2,5,1)$ at $B=20$ (Proposition~\ref{prop:rhat_rigidity}). Rows are generated by \texttt{scripts/exp\_mass\_depth\_leave\_one\_out.py}.}
\label{tab:mass_depth_leave_one_out_summary}
\end{table}

\paragraph{Audit context.}
Candidate-domain size, best/second-best mismatch context, and distribution quantiles for the mass-depth closure are recorded in Appendix~\ref{app:generated_tables} (Tables~\ref{tab:audit_closure_metrics}--\ref{tab:audit_closure_quantiles}); uncertainty-robustness stress tests are recorded in Table~\ref{tab:audit_uncertainty_robustness}, and counterfactual baseline comparisons in Table~\ref{tab:audit_counterfactual}.

\subsection{A minimal matching layer: quantized depth shifts}
\label{subsec:mass_matching_layer}

In effective field theory, scheme and threshold conventions enter as multiplicative matching factors.
In the resolution coordinate, a matching factor is an additive shift:
if $\mu_{\mathrm{ref}}$ is a reference mass and $\mu_{\mathrm{pred}}$ is a closed template value, then
\[
\Delta r:=r(\mu_{\mathrm{ref}})-\widehat r=\log_\varphi\!\left(\frac{\mu_{\mathrm{ref}}}{\mu_{\mathrm{pred}}}\right),
\qquad
\frac{\mu_{\mathrm{ref}}}{\mu_{\mathrm{pred}}}=\varphi^{\Delta r}.
\]

To make the matching layer auditable at bounded complexity, we record a minimal dyadic quantization convention:
matching shifts are summarized on the quarter-depth lattice $\Delta r\approx k/4$.
Table~\ref{tab:mass_matching_layer} reports $\Delta r$, the nearest $k/4$, the residual $\Delta r-k/4$, and the implied quantized matching factor $\varphi^{k/4}$.
Table~\ref{tab:mass_matching_layer_summary} summarizes the residual sizes across the same rows.
This does not change the closed depth formula; it makes the scheme/matching correction explicit in the same discrete language as the rest of the paper.

\begin{remark}[Why a quarter-depth lattice is the minimal dyadic choice]\label{rem:quarter_depth_minimal_dyadic}
Two structural features of the present paper single out dyadic denominators as the natural low-complexity family for matching-layer summaries.
First, phases are treated as dyadic registers $\ZZ_{2^p}$ and the holonomy diagnostics treat denominators of the form $\mathrm{denom}=2^p$ in CAP audit form (Remark~\ref{rem:dyadic_phase_register}); thus dyadic refinement is the protocol-native notion of ``one more bit'' of resolution.
Second, several candidate families used for mixing amplitudes are square-root structured (e.g.\ $\varphi^{-k/2}$ and $1/\sqrt{d}$ in \eqref{eq:ckm_candidate_family} and the analogous PMNS family), so half-depth exponents appear already at the amplitude level.
When such half-depth normalizations are composed with a logarithmic matching dictionary (as in \eqref{eq:r_of_mu_z128}), quarter-depth steps arise as the smallest dyadic lattice that can express the resulting corrections without introducing higher-denominator freedom.
Accordingly, $k/4$ is used here as the \emph{minimal} dyadic reporting lattice; finer dyadic lattices $k/2^s$ ($s>2$) can be adopted at higher complexity if one wishes to audit smaller residual structure.
We emphasize that we do not optimize the lattice denominator in this paper: $1/4$ is chosen as the minimal dyadic reporting convention, and the residuals in Tables~\ref{tab:mass_matching_layer}--\ref{tab:mass_matching_layer_summary} make the remaining mismatch explicit.
\end{remark}

\begin{table}[t]
\centering
\small
\setlength{\tabcolsep}{6pt}
\renewcommand{\arraystretch}{1.15}
\begin{tabular}{lrrrrrr}
\toprule
field & $r(\mu)$ & $\widehat r$ & $\Delta r$ & nearest & residual & $\varphi^{k/4}$ \\
\midrule
$e$ & 0.000 & 0 & +0.000 & $0/4$ & +0.000 & $1$ \\
$\mu$ & 11.080 & 11 & +0.080 & $0/4$ & +0.080 & $1$ \\
$\tau$ & 16.945 & 17 & -0.055 & $0/4$ & -0.055 & $1$ \\
$u$ & 2.996 & 6 & -3.004 & $-12/4$ & -0.004 & $0.236068$ \\
$d$ & 4.598 & 5 & -0.402 & $-2/4$ & +0.098 & $0.786151$ \\
$s$ & 10.814 & 12 & -1.186 & $-5/4$ & +0.064 & $0.547981$ \\
$c$ & 16.247 & 12 & +4.247 & $17/4$ & -0.003 & $7.73032$ \\
$b$ & 18.722 & 23 & -4.278 & $-17/4$ & -0.028 & $0.129361$ \\
$t$ & 26.456 & 28 & -1.544 & $-6/4$ & -0.044 & $0.485868$ \\
$W$ & 24.866 & 25 & -0.134 & $-1/4$ & +0.116 & $0.886652$ \\
$Z$ & 25.128 & 25 & +0.128 & $1/4$ & -0.122 & $1.12784$ \\
$H$ & 25.788 & 26 & -0.212 & $-1/4$ & +0.038 & $0.886652$ \\
\bottomrule%
\end{tabular}
\caption{Minimal matching-layer summary in the resolution coordinate. The ``nearest'' column reports the closest quarter-step $k/4$ to the observed $\Delta r$. Rows are generated by \texttt{scripts/exp\_mass\_matching\_layer.py}.}
\label{tab:mass_matching_layer}
\end{table}

\begin{table}[t]
\centering
\small
\setlength{\tabcolsep}{10pt}
\renewcommand{\arraystretch}{1.15}
\begin{tabular}{lr}
\toprule
statistic & value \\
\midrule
entries & 12 \\
median $|\Delta r-k/4|$ & 0.050 \\
p90 $|\Delta r-k/4|$ & 0.116 \\
max $|\Delta r-k/4|$ & 0.122 \\
$N_{|\cdot|\le 0.01}$ & 3 \\
$N_{|\cdot|\le 0.05}$ & 6 \\
\bottomrule%
\end{tabular}
\caption{Residual-size summary for the quarter-step matching lattice in Table~\ref{tab:mass_matching_layer}. Rows are generated by \texttt{scripts/exp\_mass\_matching\_layer.py}.}
\label{tab:mass_matching_layer_summary}
\end{table}


